\makeglossaries

\newglossaryentry{ATS}
{
    name={ATS (Automatic Text Simplification)},
    description={Automatische maschinelle Textvereinfachung zur gezielten Reduktion linguistischer Komplexität auf syntaktischer, lexikalischer und struktureller Ebene.}
}
\newglossaryentry{LS}
{
    name={LS (Leichte Sprache)},
    description={Stark reglementierte und gesetzlich normierte barrierefreie Sprachvariante des Deutschen zur Maximierung der Verständlichkeit für Menschen mit Lernschwierigkeiten oder kognitiven Einschränkungen.}
}
\newglossaryentry{AS}
{
    name={AS (Alltagssprache)},
    description={Deutsche Standardsprache und behördliche Fachsprache der ursprünglichen Ausgangstexte vor der Vereinfachung.}
}
\newglossaryentry{DPO}
{
    name={DPO (Direct Preference Optimization)},
    description={Verfahren zur mathematisch geschlossenen Präferenzoptimierung generativer Sprachmodelle auf Paarvergleichen ohne separates Belohnungsmodell \citep{rafailov2023direct}.}
}
\newglossaryentry{SFT}
{
    name={SFT (Supervised Fine-Tuning)},
    description={Überwachtes Feintuning neuronaler Sprachmodelle auf parallelen Textpaaren zur Etablierung grundlegender Übersetzungsmuster.}
}
\newglossaryentry{LoRA}
{
    name={LoRA (Low-Rank Adaptation)},
    description={Parametereffiziente Feintuning-Methode, welche die Gewichtsmatrizen neuronaler Schichten durch niederrangige Matrizenzerlegungen adaptiert \citep{hu2022lora}.}
}
\newglossaryentry{SARI}
{
    name={SARI},
    description={Referenzbasierte Evaluierungsmetrik für Textvereinfachungen, die N-Gramm-Präzision und -Recall für Hinzufügungen, Löschungen und Beibehaltungen separat quantifiziert \citep{xu2016optimizing}.}
}
\newglossaryentry{MATTR}
{
    name={MATTR (Moving-Average Type-Token Ratio)},
    description={Längenunabhängiges Maß der lexikalischen Vielfalt durch Berechnung der Type-Token-Ratio über ein gleitendes Wortfenster \citep{covington2010cutting}.}
}
\newglossaryentry{BITV}
{
    name={BITV 2.0},
    description={Barrierefreie-Informationstechnik-Verordnung; deutsche Rechtsverordnung zur barrierefreien Gestaltung digitaler Angebote öffentlicher Stellen nach dem BGG.}
}